%% =====================================================================
%%  How to use this book : la structure d'un chapitre et les repères
%% =====================================================================
\liminaire{mode}{How to use this book}{A chapter in seven steps, always in the same order}

The twenty chapters of this book are all built on the same model. After an opening page that
presents the big picture, your goals and your study plan, each chapter goes through seven kinds of
sections, which you can recognise by their colour and their icon.\par\vspace{5mm}

\newcommand{\otEtape}[4]{% clé de rubrique, titre, texte, durée
  \noindent\begin{tikzpicture}[baseline=(t.north)]
    \fill[\otRubCoul{#1},rounded corners=2.4mm] (0,-11mm) rectangle (11mm,0);
    \otDessinIcone{\otRubIco{#1}}{white}{\otRubCoul{#1}}{6.6mm}{(5.5mm,-5.5mm)}
    \node[anchor=north west,inner sep=0pt] (t) at (0,0) {};
  \end{tikzpicture}\hspace{4mm}%
  \begin{minipage}[t]{\dimexpr\linewidth-15mm\relax}
    {\otSemi\fontsize{7.2}{9}\selectfont\color{\otRubCoul{#1}}\otCaps[12]{\otRubNom{#1}}}\hspace{2mm}%
    {\otSemi\fontsize{12}{14}\selectfont #2}\hfill{\otRegular\fontsize{8.4}{10}\selectfont\color{otMuted}#4}\par\vspace{1mm}
    {\fontsize{10.2}{13}\selectfont #3\par}
  \end{minipage}\par\vspace{4.6mm}}

\otEtape{learn}{The lessons}{New ideas, one at a time, each with short programs that you run
  and modify. Little \enquote{Your turn} tasks let you try straight away.}{90 min}
\otEtape{hood}{How Python works}{What really happens inside the computer when your program runs.
  Understanding the machine makes everything else easier.}{15 min}
\otEtape{debug}{Errors and how to fix them}{The mistakes that everybody makes with the new ideas
  of the chapter, the messages Python shows, and how to correct them.}{15 min}
\otEtape{pro}{Professional habits}{The way experienced programmers write and organise their code,
  from your very first programs.}{10 min}
\otEtape{build}{A mini-project}{A small complete program, built step by step, that uses everything
  in the chapter.}{20 min}
\otEtape{practice}{Exercises}{Twelve progressive exercises: predict, fix, complete, write\ldots{} All
  are solved and explained at the end of the book.}{45 min}
\otEtape{review}{Chapter review}{The essentials on one page, a ten-question quiz and a checklist to
  see what you have mastered.}{15 min}

After every four chapters, a \textbf{guided project} and a \textbf{checkpoint} help you combine and
consolidate what you have learnt.

\clearpage
\otIntertitre{Code, output and errors}
\begin{code}{file=example.py}
# A program file: type it in the IDLE editor
print("Hello!")
\end{code}
\begin{sortie}
Hello!
\end{sortie}
\begin{tcolorbox}[otBase,colback=white,colframe=otRule,boxrule=0.5pt]
\begin{tabularx}{\linewidth}{@{}>{\RaggedRight\arraybackslash}p{30mm}@{\hspace{5mm}}Y@{}}
\textbf{Light box} & A \textbf{program} that you type in the editor. Its file name is at the top;
  the line numbers in the margin are not part of the code.\\[2.4mm]
\textbf{Dark box: Output} & What the program displays when you run it. Text that the user types on
  the keyboard is shown in {\color{otGreenDark}\ttfamily\bfseries green}.\\[2.4mm]
\textbf{Dark box: Python shell} & A conversation with Python in the IDLE Shell. Type only what comes
  after the prompt \py{>>>}.\\[2.4mm]
\textbf{Red edge: Error message} & What Python displays when something goes wrong. Read it from the
  bottom up.\\[2.4mm]
\otRepere{1}\hspace{1mm}\otRepere{2} & Numbered markers link a line of code to its explanation
  just below.\\
\end{tabularx}
\end{tcolorbox}

\otIntertitre{The boxes}
\begin{paire}
\begin{regle}[Key idea]
The rule or idea to remember.
\end{regle}
\begin{attention}[Watch out]
A trap to avoid:\par\vspace{1mm}
\fauxjuste{print "Hi"}{print("Hi")}
\end{attention}
\end{paire}
\begin{paire}
\begin{astuce}[Tip]
Practical advice to learn faster.
\end{astuce}
\begin{pratique}[Pro habit]
How professionals do it.
\end{pratique}
\end{paire}
\begin{paire}
\begin{culture}[Did you know?]
History, anecdotes and the world of programming.
\end{culture}
\begin{definition}{key term}
A word of the programmer's vocabulary, defined simply. All of them are in the glossary.
\end{definition}
\end{paire}

\otIntertitre{The exercises}
\begin{tcolorbox}[otBase,colback=white,colframe=otRule,boxrule=0.5pt]
\begin{tabularx}{\linewidth}{@{}p{34mm}Y@{}}
\otBadgeType{predict} & say what a program displays, without running it\\[1.6mm]
\otBadgeType{fix} & find and correct the mistakes in a program\\[1.6mm]
\otBadgeType{complete} & fill in the blanks of a program\\[1.6mm]
\otBadgeType{write} & write a program from scratch\\[1.6mm]
\otBadgeType{shell} & experiment in the Python shell\\[1.6mm]
\otBadgeType{read} & read some code or an error message and answer questions\\[1.6mm]
\otBadgeType{explain} & explain in your own words\\[1.6mm]
\otBadgeType{challenge} & a harder problem, to go further\\[2.6mm]
\otDifficulte{1}\hspace{2mm}\otDifficulte{3} & the difficulty, from one dot (easy) to three dots (demanding)\\[1.6mm]
{\otLight\fontsize{7.5}{9}\selectfont\color{otMuted}solution p.~123} & where to find the solution, at the end of the book\\
\end{tabularx}
\end{tcolorbox}

\otIntertitre{At the end of the book}
\begin{tcbraster}[raster columns=3,raster equal height=rows,raster column skip=3.5mm,raster before skip=1mm]
\otCartePrg{coche}{otGreenDark}{Solutions}{Every exercise solved, with the reasoning, the complete code and its output, and the classic mistakes.}
\otCartePrg{livre}{otNavy}{Glossary}{Every key term of the volume with a short definition and the chapter where it is explained.}
\otCartePrg{telecharger}{otOrangeDark}{Code online}{All the programs of the book, ready to run, to download from the website.}
\end{tcbraster}
